% 对应题号 M1-47，复用已核对的本题几何图。
% 修订：启用 xeCJK 自动换行；保留原题参数与物理关系。
\documentclass[tikz,border=14pt]{standalone}
\usepackage{fontspec}
\usepackage{xeCJK}
\usepackage{amsmath,amssymb}
\usepackage{xcolor}
\usepackage{tikz}
\usetikzlibrary{arrows.meta,calc,patterns,positioning,bending}
\IfFontExistsTF{Heiti SC}{\IfFontExistsTF{Heiti SC}{\setmainfont{Heiti SC}}{\setmainfont{FandolHei-Regular}}\IfFontExistsTF{Heiti SC}{\setsansfont{Heiti SC}}{\setsansfont{TeX Gyre Heros}}}{\IfFontExistsTF{Songti SC}{\setmainfont{Songti SC}}{\setmainfont{FandolSong-Regular}}\IfFontExistsTF{Songti SC}{\setsansfont{Songti SC}}{\setsansfont{FandolSong-Regular}}}
\definecolor{bluefield}{RGB}{24,76,255}
\definecolor{redfield}{RGB}{232,49,35}
\definecolor{orangefield}{RGB}{236,132,18}
\definecolor{greenfield}{RGB}{18,132,86}
\definecolor{purplefield}{RGB}{112,72,190}
\definecolor{linegray}{RGB}{60,68,78}
\definecolor{softblue}{RGB}{238,243,255}
\definecolor{softgreen}{RGB}{238,249,244}
\definecolor{softorange}{RGB}{255,247,235}
\tikzset{
  arrow/.style={-{Latex[length=3.1mm,width=2.1mm]},line width=1.1pt},
  thinarr/.style={-{Latex[length=2.5mm,width=1.7mm]},line width=.85pt},
  axis/.style={-{Latex[length=2.5mm,width=1.7mm]},line width=.8pt,draw=linegray},
  guide/.style={draw=linegray!55,dashed,line width=.65pt},
  box/.style={draw=linegray,rounded corners=3pt,fill=white,align=center,inner sep=6pt,line width=.75pt},
  bluebox/.style={box,draw=bluefield,fill=softblue},
  greenbox/.style={box,draw=greenfield,fill=softgreen},
  orangebox/.style={box,draw=orangefield,fill=softorange}
}
\IfFontExistsTF{Heiti SC}{\setCJKmainfont[AutoFakeBold=1.5]{Heiti SC}}{\setCJKmainfont[AutoFakeBold=1.5]{FandolHei-Regular}}
\IfFontExistsTF{Heiti SC}{\setCJKsansfont{Heiti SC}}{\setCJKsansfont{FandolHei-Regular}}
\begin{document}

\begin{tikzpicture}[x=1cm,y=1cm]
\path[use as bounding box] (-10,-4.6) rectangle (10,4.6);
\fill[white] (-10,-4.6) rectangle (10,4.6);
\node[font=\Large\bfseries,text width=18.6cm,align=center] at (0,4.05) {刚体转动速度场：旋度沿转轴，大小是角速度的两倍};
\begin{scope}[shift={(-4.7,-.2)}]
\draw[linegray] (0,0) circle(2.25);
\foreach \a in {0,45,...,315}{
  \draw[arrow,bluefield] ({1.55*cos(\a)},{1.55*sin(\a)})--({1.55*cos(\a)-.75*sin(\a)},{1.55*sin(\a)+.75*cos(\a)});
}
\node[redfield,font=\Large] at (0,0) {$\odot$};
\node[redfield] at (0,.55) {$\boldsymbol\omega$};
\node[bluefield] at (0,-2.8) {$\mathbf v=\boldsymbol\omega\times\mathbf r$};
\end{scope}
\begin{scope}[shift={(4.6,-.1)}]
\node[bluebox,text width=6.7cm] at (0,1.65) {相对圆心两侧的速度差同时贡献局部环量};
\node[greenbox,text width=6.7cm] at (0,0) {$\displaystyle \nabla\times\mathbf v=2\boldsymbol\omega$};
\node[orangebox,text width=6.7cm] at (0,-1.7) {$\displaystyle \boldsymbol\omega_{\rm 局部}=\frac12\nabla\times\mathbf v$\\系数 $2$ 来自速度梯度的反对称部分。};
\end{scope}
\node[font=\small,text width=18cm,align=center] at (0,-4.1) {从角速度的正向看，图示绕轴逆时针旋转。};
\end{tikzpicture}
\end{document}
